\section{Результаты}
\subsection{Последовательный контрольный процесс выделяет эффект отбора компонентов}
В основном routing-сценарии P0 retained-load снижает среднюю стоимость на 38.2\% относительно static planning и на 38.1\% относительно capacity-at-20 (табл.~\ref{tab:primary}). Парные 97.5\%-ные интервалы с поправкой на множественные сравнения равны $[-51.56,-19.40]$ и $[-51.23,-19.32]$ тыс. Относительно predictive selection снижение составляет 35.4\%, с описательным 95\%-ным интервалом $[-46.30,-16.66]$ тыс.

\begin{table}[t]\caption{P0: средняя реализованная стоимость в тысячах по 12 независимым обучениям на механизм. Все методы используют одинаковые тестовые условия.}\label{tab:primary}\centering\small\setlength{\tabcolsep}{3pt}
\begin{tabular}{lrrrrrr}\toprule
Механизм & \shortstack{Без\\действия} & Static & Retained & Predictive & \shortstack{Capacity\\at 20} & \shortstack{Rework\\at 20}\\\midrule
Фиксированный & 23.6 & 18.0 & 19.9 & 18.0 & 17.7 & 22.6 \\
Обслуживание & 211.7 & 111.0 & 111.0 & 111.0 & 111.0 & 202.7 \\
Возвраты & 147.8 & 92.8 & 57.4 & 88.8 & 92.6 & 45.3 \\
Совместный & 375.2 & 309.2 & 288.3 & 292.5 & 309.2 & 284.1 \\
\bottomrule

\end{tabular}\end{table}

Предиктивный селектор сохраняет нагрузочный компонент возвратов в 2 из 12 обучений и выбирает сокращение rework в 1; retained-load выбирает его в 9. Все варианты retained-load используют порог 20, поэтому результат относится к выбору действия, а не улучшению времени активации. Фиксированное сокращение rework и планировщик с известным генератором стоят по 45.3 тыс., меньше 57.4 тыс. у retained-load. Парная разность равна $12.04$ тыс., 95\%-ный интервал — $[-1.87,25.94]$. Эвристика силы обратной связи стоит 73.0 тыс. и выбирает rework в 5 случаях; интервал разности retained-load и эвристики равен $[-38.90,7.65]$ тыс. Эти сравнения не устанавливают преимущества над сильными простыми альтернативами.

Планировщик с известным генератором выбирает rework-at-20 во всех 12 случаях. Превышение стоимости retained-load связано с 3 выборами capacity: обученные модели меняют порядок предпочтения действий относительно независимого теста. В этих обучающих историях WIP достигает лишь 12--19 экземпляров. Ошибка относится к ранжированию действий за пределами наблюдавшейся области; эксперимент не выделяет единственный ошибочный коэффициент.

Другие механизмы отделяют полезную адаптацию от безусловного усложнения модели. При service feedback все планировщики выбирают capacity-at-20. При fixed retained-load повышает стоимость на 10.9\% относительно static planning; при combined снижает её на 6.7\%, но частота тяжёлых событий растёт с 69.3\% до 71.0\%. Снижение интегральной стоимости не обязательно уменьшает число длительных перегрузок.

\subsection{Чувствительность селектора и наблюдаемая реакция}
Диагностика P0 повторно использует все сохранённые тестовые действия, изменяя порог отбора возвратов по сетке $0,0.001,0.0025,0.005,0.01$. Компонент сохраняется соответственно в $8,5,2,2,0$ из 12 routing-обучений; средние стоимости составляют $80.8,84.9,88.8,88.8,92.8$ тыс. Стоимость retained-load остаётся $57.4$ тыс. Следовательно, наблюдаемое различие не ограничивается исходным порогом $0.005$. При порогах отбора длительности $0,0.0025,0.005,0.01,0.02$ стоимость селектора равна $85.1$ тыс. при нулевом пороге и $88.8$ при остальных. Это диагностика на прежнем независимом тестовом каталоге, а не новые подтверждающие наборы. Сохранение только зависимости длительности стоит $69.4$ тыс., только возвратов — $80.8$ тыс.; оба компонента дают минимальную среднюю стоимость среди этих абляций.

Во всех 289 routing-траекториях P0 с тяжёлым событием без вмешательства правило retained-load срабатывает заранее, со средним опережением 108.0. Оно предотвращает 93 события (32.2\%). В 192 спокойных траекториях активаций нет, однако правило срабатывает в 95 из 287 траекторий без тяжёлого события в варианте no action. При неизменных исходных политиках фактическое сокращение rework на 30\%, 50\% и 70\% даёт снижение стоимости на 17.0--48.7\% относительно static planning. Эти результаты разделяют своевременность активации, предотвращение события, снижение стоимости и устойчивость к ошибке в эффекте действия.

\subsection{Как структура меняет пользу сохранения нагрузочных зависимостей}
Все 288 новых обучений завершены, все 3168 калибровочных/валидационных эпизодов удовлетворяют предкризисному условию. Таблица~\ref{tab:structures_results} сравнивает решения при нагрузочных возвратах и заданной структуре. Retained-load снижает среднюю стоимость относительно предиктивного отбора на 8.1\% в P1, 29.8\% в P2, 34.6\% в P5 и 26.4\% в P6. Стоимость возрастает на 3.0\% в P3 и 15.7\% в P4. Ни одно из четырёх улучшений не исключает ноль после поправки на шесть основных сравнений; рост P4 исключает его: скорректированный интервал разности равен $[1.9,8.3]$ тысяч единиц стоимости. Поэтому расширение не устанавливает подтверждающего улучшения для отдельной структуры.

\begin{table}[t]\caption{Routing feedback при известной структуре: стоимость каждого обученного варианта нормирована на соответствующую стоимость без вмешательства. Regret равен разности стоимости retained-load и планировщика с известным генератором, также нормированной на no action. Каждая строка содержит 12 независимых обучений.}\label{tab:structures_results}\centering\small\setlength{\tabcolsep}{4pt}
\begin{tabular}{lrrrrrr}\toprule
Процесс & Static & Predictive & Retained & Heuristic & Generator & Regret\\\midrule
P1 & 0.587 & 0.549 & 0.502 & 0.501 & 0.429 & 0.073 \\
P2 & 0.731 & 0.614 & 0.428 & 0.514 & 0.428 & 0.000 \\
P3 & 0.630 & 0.630 & 0.650 & 0.630 & 0.516 & 0.133 \\
P4 & 0.520 & 0.528 & 0.610 & 0.620 & 0.512 & 0.098 \\
P5 & 1.030 & 0.727 & 0.473 & 0.422 & 0.422 & 0.050 \\
P6 & 0.575 & 0.575 & 0.422 & 0.460 & 0.349 & 0.073 \\
\bottomrule\end{tabular}\end{table}


Вторичное сравнение с равными весами показывает, где возникает средняя польза. При нагрузочных возвратах нормированная разность retained минус predictive равна $-0.090$ с описательным 95\%-м интервалом $[-0.135,-0.045]$; при обеих обратных связях — $-0.036$ $[-0.052,-0.020]$. При отсутствии связей и только нагрузочной длительности она, напротив, возрастает на 0.058 и 0.069. Это разности стоимости, делённой на стоимость без вмешательства для каждого обучения, а не проценты относительно предиктивного планировщика. Результат поддерживает условную пользу сохранения зависимостей в исследованном семействе, а не правило включать все нагрузочные компоненты.

\begin{figure}[t]\centering
\includegraphics[width=\textwidth]{generated/structural_matrix_ru.pdf}
\caption{Снижение средней стоимости retained-load относительно предиктивного отбора, в процентах. Положительные значения означают пользу сохранения зависимостей. Панели разделяют заданную и обнаруженную структуры. Звёздочка отмечает ненулевое основное сравнение после поправки на шесть проверок; остальные ячейки описательные.}\label{fig:structural_matrix}
\end{figure}

\subsection{Ошибки обнаружения структуры могут менять подготовленное действие}
При обнаруженной структуре общий результат для нагрузочных возвратов становится $0.001$ $[-0.037,0.038]$: преимущество известной структуры не переносится на полный конвейер. P2 показывает последствия для решения. При заданной структуре retained-load во всех 12 обучениях выбирает сокращение возвратов, совпадающее с действием планировщика по генератору, и получает нормированную стоимость 0.428. При обнаруженной структуре он во всех обучениях выбирает увеличение мощности и получает 0.897. Обучающие данные, поведенческие параметры, ресурсы и набор действий между вариантами неизменны. Следовательно, структурное представление меняет вмешательство, которое поддерживают одни и те же поведенческие оценки. P1, P2 и P6 добавляют пути во всех обнаруженных моделях при fitness 1.0; средняя precision равна 0.497, 0.642 и 0.487. P3, P4 и P5 восстанавливают эквивалентные видимые языки и повторяют политики заданной структуры во всех 48 обучениях каждого процесса. Поэтому соответствие наблюдавшимся трассам само по себе не устанавливает корректность решений.

\subsection{Что симуляция добавляет к локальной эвристике}
P3 и P5 предлагают три содержательно различных вмешательства, но сравнения не устанавливают преимущества retained-load перед эвристикой силы обратной связи. В P5 с нагрузочными возвратами правило всегда выбирает rework-at-10 и практически совпадает с планированием по генератору; retained-load выбирает этот тип действия в 11 из 12 обучений. В P3 планирование по генератору всегда воздействует на локальный возврат, а retained-load — лишь в 1 из 12 обучений при точной структуре. Разрыв возникает в обученном ранжировании действий. В P4 удаление нагрузочных членов маршрутизации снижает нормированную стоимость с 0.610 до 0.516; удаление нагрузочных членов длительности оставляет её на уровне 0.611. Абляция локализует ошибку в оценённых зависимостях возвратов, не выбирая исправленный метод по тестовым исходам.

Повтор P5 показывает взаимодействие ресурсов. При обеих обратных связях дополнительный работник B уменьшает его среднюю очередь с 141.02 до 40.19, но увеличивает очередь D с 1.49 до 118.56. Наибольшая очередь перемещается от B к D во всех 24 парных диагностических траекториях. Ускорение ранней стадии направляет работу на следующую стадию, где длительность и повторения тоже зависят от нагрузки. Очереди измерены на исходном тестовом горизонте при фиксированных правилах действий.

P6 с нагрузочной длительностью даёт иной диагностический результат. Локальное правило обычно добавляет мощность параллельной ветви; имитационный retained-load чаще выбирает общий пул и получает нормированную стоимость 0.633 против 0.920. Описательный парный интервал разности стоимости равен $[-26.3,-7.0]$ тысяч. У статического планирования средняя стоимость ниже: 0.494, близко к 0.492 у генератора. Сравнение поддерживает оценку системных последствий, когда локальные показатели вводят в заблуждение, и одновременно показывает, что дополнительные нагрузочные компоненты не нужны для этого конкретного улучшения.

\subsection{Чувствительность к отбору компонентов и эффектам действий}
Перебор селектора сохраняет направление сравнения при нагрузочных возвратах в P1, P2, P5 и P6: retained-load имеет меньшую среднюю стоимость при каждой проверенной настройке любого из двух порогов. Например, изменение $\Delta_R$ от 0 до 0.01 меняет нормированную стоимость предиктивного отбора в P5 с 0.575 до 0.879 против 0.473 у retained-load. В P4 предиктивный отбор лучше во всём диапазоне; в P3 порядок меняется вместе с порогом возвратов. Следовательно, эффект шире одной настройки по умолчанию, но зависит от процесса. Это диагностические сравнения без выбора нового порога по тесту.

При отдельном изменении состава P2 до 65\% сложных экземпляров неизменные политики retained и predictive получают нормированную стоимость 0.623 и 0.699. Неопределённость вмешательства существеннее в P3: при фактическом сокращении возвратов на 30\% и 70\% retained-load стоит 0.742 и 0.573 против 0.671 и 0.600 у предиктивного отбора. Порядок меняется при фиксированных политиках. Поэтому карточка реагирования должна указывать диапазон эффектов, в котором рекомендация была проверена.

